\section{Graphs, Neighbors, and Networks}

We have been talking about cells, neighbors, and touching. There is a clean, general way to think about all of it, and it strips the crystal down to its essentials. It is the idea of a graph. Do not let the word scare you, a graph here just means dots and connections.

\subsection{Dots and lines}

A graph is two things: a set of dots, and a set of lines joining some of the dots. That is all. The dots are usually called nodes, the lines are called edges. A subway map is a graph: stations are dots, tracks are lines. A friendship network is a graph: people are dots, friendships are lines.

For us, the crystal is a graph. Each cell is a dot. Each shared edge between two touching cells is a line. When we strip away the pretty hyperbolic drawing and the cell shapes, what is left, what the theory actually runs on, is a graph: a huge collection of dots, each joined to a handful of others.

\subsection{Only the connections matter}

Here is the crucial mental shift. Once you see the crystal as a graph, you realize the picture, the positions, the cell shapes, are not the real content. The real content is who is connected to whom. Two crystals drawn completely differently are the same crystal if their connection patterns match. You can stretch the drawing, slide it, redraw it in a different style, and as long as the dots keep the same neighbors, nothing physical has changed.

This is why earlier we said the Poincare disk is a map of relations, not a container. The space is not a stage the cells sit on. The space is just a way of visualizing the connections. The connections are primary.

\subsection{Degree: how many neighbors}

One number describes a dot: its degree, the count of lines coming out of it, that is, how many neighbors it has. In the crystal each cell has a small, roughly fixed degree, a handful of neighbors, because each cell touches only the cells right next to it. This ``small and local'' property is important. A cell never talks to far-away cells directly. It only ever touches its immediate neighbors. All long-range effects have to be built up step by step, neighbor to neighbor.

That locality is not a limitation to apologize for. It is exactly the feature that gives the model a speed limit, a maximum rate at which influence can spread, which becomes the speed of light later on. Because every cell only touches its neighbors, nothing can jump across the crystal instantly. It has to travel, cell by cell.

\subsection{Distance as steps}

On a graph, distance is just counting steps. The distance between two dots is the fewest lines you have to cross to get from one to the other. This step-counting distance is what ``space'' will mean for us. Two cells are near if you can get between them in few steps, far if it takes many. There is no ruler floating above the crystal. Distance is woven from the connections themselves.

So when later chapters say ``space emerges,'' this is the seed of what they mean. Space is not assumed. It is the step-counting distance on a graph, the shadow cast by the pattern of who touches whom.

\subsection{The whole model, in graph terms}

We can now state the whole substrate in one plain sentence. It is a graph: a vast set of dots, each joined to a few neighbors, arranged so that step-counting distance behaves like curved three-dimensional space, growing new dots at its frontier forever. Everything else in the walkthrough, the tones, the rule, matter, mind, lives on top of this graph. Hold the graph in mind as the bare skeleton, and the rest as what happens on it.

\textbf{Takeaway:} a graph is just dots joined by lines, and the crystal is one: cells are dots, shared edges are lines. Only the connections matter, not the drawing. Each cell touches just a few neighbors, which gives a speed limit, and distance is simply the number of steps between dots, which is what space will turn out to be.
